\LectureAnswerSection

\begin{drillsolution}{n3d1}\textbf{在电梯系中求截面以下的有效重量}\par\nobreak
图\ref{n3d1q}电梯系有效重力 $g_e=g+a_0$ 向下。自然坐标从顶端向下取，截面以下总质量为 $M+m_0(1-u/L_0)$，所以 $T(u)=[M+m_0(1-u/L_0)]g_e$。
\fig{n3d1s}{末端载荷使整个轴力图提高同一个常量}{\begin{tikzpicture}
\draw[axis](0,0)--(4.6,0)node[right]{$u$};\draw[axis](0,0)--(0,2.7)node[above]{$T$};\draw[world](0,2.2)--(4,.8);\draw[guide](0,.8)--(4,.8);\node[left]at(0,2.2){$(M+m_0)g_e$};\node[right]at(4,.8){$Mg_e$};\node[below]at(4,0){$L_0$};
\end{tikzpicture}}


\par\medskip\Needspace{4\baselineskip}\noindent\textbf{累加局部伸长并检查失重情形}\par\nobreak
按图\ref{n3d1s}，自然微元 $\dd u$ 的局部劲度为 $kL_0/\dd u$，伸长为 $T(u)\dd u/(kL_0)$。从上端积到下端：
\begin{align*}
\Delta L&=\frac{g+a_0}{kL_0}\int_0^{L_0}
\left[M+m_0-\frac{m_0u}{L_0}\right]\dd u\\
&=\frac{g+a_0}{kL_0}
\left[(M+m_0)u-\frac{m_0u^2}{2L_0}\right]_0^{L_0}\\
&=\frac{g+a_0}{kL_0}
 \left[(M+m_0)L_0-\frac{m_0L_0}{2}\right]
=\frac{(M+m_0/2)(g+a_0)}k.
\end{align*}
上端力为 $(M+m_0)(g+a_0)$，下端轴力为 $M(g+a_0)$。若电梯向下加速度恰为 $g$，有效重力为零，原来已相对静止且无振动的极限状态可无拉伸；从有伸长状态突然自由落体则会发生弹性振动，不能直接假设形变瞬间归零。
\end{drillsolution}

\begin{drillsolution}{n3d2}\textbf{由径向压差求两液面高度差}\par\nobreak
图\ref{n3d2q}中径向微元满足 $\dd p=\rho\omega^2r\,\dd r$，同高底管两端的压差为
\[
p(r_2)-p(r_1)=\rho\omega^2\int_{r_1}^{r_2}r\,\dd r
=\frac12\rho\omega^2[r^2]_{r_1}^{r_2}
=\frac12\rho\omega^2(r_2^2-r_1^2).
\]
各支管从液面向下积分给 $p(r_i)-p_0=-\rho g[z]_{h_i}^0=\rho gh_i$，两式相减，故 $\delta h=h_2-h_1=\omega^2(r_2^2-r_1^2)/(2g)$。
\fig{n3d2s}{等体积转移使窄管下降多、宽管上升少}{\begin{tikzpicture}
\fill[blue!10](-1.5,0)rectangle(2.2,.2);\fill[blue!10](-1.5,0)rectangle(-.8,.7);\fill[blue!10](1,0)rectangle(2.2,1.65);\fill[red!10](-1.5,.7)rectangle(-.8,1.3);\fill[blue!25](1,1.3)rectangle(2.2,1.65);\draw[world](-1.5,2.5)--(-1.5,0)--(2.2,0)--(2.2,2.5);\draw[world](-.8,2.5)--(-.8,.2)--(1,.2)--(1,2.5);\draw[guide](-2,1.3)--(2.8,1.3)node[right]{$h$};\draw[blue,thick](-1.5,.7)--(-.8,.7) (1,1.65)--(2.2,1.65);\node[left]at(-1.5,.7){$h_1$};\node[right]at(2.2,1.65){$h_2$};\node at(.35,-.6){$S_1(h-h_1)=S_2(h_2-h)$};
\end{tikzpicture}}


\par\medskip\Needspace{4\baselineskip}\noindent\textbf{用等体积交换确定各液面}\par\nobreak
图\ref{n3d2s}中 $S_1h_1+S_2h_2=(S_1+S_2)h$，代入 $h_2=h_1+\delta h$，解得 $h_1=h-S_2\delta h/(S_1+S_2)$、$h_2=h+S_1\delta h/(S_1+S_2)$。从轴积分到内侧支管底部，$p(r_1)-p(0)=\rho\omega^2[r^2/2]_0^{r_1}$，故轴处压强为
\begin{align*}
 p&=p_0+\rho gh_1-\tfrac12\rho\omega^2r_1^2\\
 &=p_0+\rho gh-\frac{\rho\omega^2}{2}
 \left[\frac{S_2(r_2^2-r_1^2)}{S_1+S_2}+r_1^2\right]\\
 &=p_0+\rho gh-\frac{\rho\omega^2}{2}
 \frac{S_1r_1^2+S_2r_2^2}{S_1+S_2}.
\end{align*}
最后一步将括号内通分，$-S_2r_1^2$与 $(S_1+S_2)r_1^2$中的 $S_2r_1^2$相消。


\par\medskip\Needspace{4\baselineskip}\noindent\textbf{筛选首先到达的边界}\par\nobreak
令 $h_1=0$ 得 $\omega_{\rm dry}^2=2gh(S_1+S_2)/[S_2(r_2^2-r_1^2)]$。若先发生外侧溢出或空化，则这一见底阈值不再实际达到，必须按先发生的边界切换模型。
\end{drillsolution}

\begin{drillsolution}{n3d3}\textbf{用离口时间表示水滴位置}\par\nobreak
图\ref{n3d3q}随桶平移系中，以离口为当地时间零点，水滴初速度为 $(0,-v_e)$、相对加速度为 $(-a,-g)$。逐分量积分得到
\begin{align*}
 v_x'(\tau)&=0+\int_0^\tau(-a)\dd s=-a\tau,\\
 v_z'(\tau)&=-v_e+\int_0^\tau(-g)\dd s=-v_e-g\tau,\\
 x'(\tau)&=\int_0^\tau(-as)\dd s
 =[-as^2/2]_0^\tau=-a\tau^2/2,\\
 z'(\tau)&=\int_0^\tau(-v_e-gs)\dd s
 =[-v_es-gs^2/2]_0^\tau=-v_e\tau-g\tau^2/2.
\end{align*}
令 $h=-z'$，整理为 $g\tau^2+2v_e\tau-2h=0$，配方得到 $(g\tau+v_e)^2=v_e^2+2gh$。由于 $\tau\ge0$，取正号分支，
\[\tau=\frac{\sqrt{v_e^2+2gh}-v_e}{g}.\]
\fig{n3d3s}{水流切向速度随离口时间改变，坐标仍取竖直向上}{\begin{tikzpicture}
\draw[axis](0,0)--(1.6,0)node[right]{$x'$};\draw[axis](0,-.2)--(0,1.3)node[above]{$z'$};\fill(0,0)circle(2pt)node[above left]{出口};\draw[world,domain=0:1.8,samples=45]plot({-.5*(\x)^2},{-.7*\x-.5*(\x)^2});\draw[flow,blue](0,0)--(0,-.65)node[right]{$v_e$};\fill(-.5,-1.2)circle(2pt);\draw[flow,blue](-.5,-1.2)--(-1,-2.05)node[left]{$\bm v_{\rm rel}$};\draw[flow,red](2,-1.2)--(1.5,-1.7)node[right]{$(-a,-g)$};\draw[guide](-1.62,-2.88)--(.5,-2.88);\draw[dim](.45,0)--(.45,-2.88)node[midway,right]{$h$};
\end{tikzpicture}}


\par\medskip\Needspace{4\baselineskip}\noindent\textbf{用切向速度换算线密度并检查极限}\par\nobreak
按图\ref{n3d3s}，$|\bm v_{\rm rel}|=\sqrt{a^2\tau^2+(v_e+g\tau)^2}$，所以 $\lambda=Q/\sqrt{a^2\tau^2+(v_e+g\tau)^2}$，其中 $\tau$ 取上式。由于 $x'$ 与 $z'$ 不再保持固定比例，$a\ne0$、$v_e>0$ 时流线为曲线。$v_e\to0$ 回到例题的斜直线与线密度；$a\to0$ 则回到竖直水流 $\lambda=Q/\sqrt{v_e^2+2gh}$。
\end{drillsolution}

\begin{drillsolution}{n3d4}\textbf{保留材料方程并改变末端受力}\par\nobreak
图\ref{n3d4q}取自然坐标 $u$，微元质量为 $(m_0/L_0)\dd u$，向外为正时
$T'(u)=-(m_0/L_0)\omega^2r$。把 $T=kL_0(r'-1)$ 求导得 $T'=kL_0r''$，所以
$r''+(\beta/L_0)^2r=0$，其中 $\beta=\omega\sqrt{m_0/k}$。
对 $\beta>0$试解 $r=e^{\lambda u}$，逐次求导为 $r'=\lambda e^{\lambda u}$、$r''=\lambda^2e^{\lambda u}$。代入微分方程并约去指数因子，得到特征方程
$\lambda^2+(\beta/L_0)^2=0$；先移项为 $\lambda^2=-\beta^2/L_0^2$，两根为 $\pm i\beta/L_0$，故实数通解为
\[
r=C\sin(\beta u/L_0)+D\cos(\beta u/L_0).
\]
固定端 $r(0)=0$ 给 $D=0$，求导得 $r'=C(\beta/L_0)\cos(\beta u/L_0)$。现在末端球要求的轴力是 $M\omega^2L$，不再为零。
\fig{n3d4s}{新边界把末端伸长梯度同末端半径联系起来}{\begin{tikzpicture}
\draw[spring](-2,0)--(1.6,0);\fill(1.6,0)circle(5pt)node[right]{$M$};\draw[force,blue](1.6,0)--(.3,0)node[above]{$T_L=M\omega^2L$};\draw[flow](1.6,.5)--(1.6,1.6)node[above]{$\omega L$};\node at(0,-.8){$T_L=kL_0[r'(L_0)-1]$};
\end{tikzpicture}}


\par\medskip\Needspace{4\baselineskip}\noindent\textbf{代入球的向心力边界求长度}\par\nobreak
由图\ref{n3d4s}，末端半径 $L=r(L_0)=C\sin\beta$，末端受力条件为
$kL_0(C\beta\cos\beta/L_0-1)=M\omega^2C\sin\beta$。把含 $C$ 的项移到同侧，
\[
C\left[k\beta\cos\beta-M\omega^2\sin\beta\right]=kL_0.
\]
由此求得 $C$，再代入 $L=C\sin\beta$ 并使用 $\omega^2/k=\beta^2/m_0$，得
\[\frac L{L_0}=\frac{\sin\beta}{\beta\cos\beta-(M/m_0)\beta^2\sin\beta}.\]


\par\medskip\Needspace{4\baselineskip}\noindent\textbf{检查正拉伸分支与质量极限}\par\nobreak
取分母首次变零之前的正拉伸分支。$M\to0$时，分母的含 $M$项消失，$L/L_0=\sin\beta/(\beta\cos\beta)=\tan\beta/\beta$。

$m_0\to0$时 $\beta\to0$，但 $(M/m_0)\beta^2=M\omega^2/k$不趋零。先将分子、分母同时除以 $\beta$，再取极限：
\[
 \frac L{L_0}
 =\frac{\sin\beta/\beta}
 {\cos\beta-(M\omega^2/k)(\sin\beta/\beta)}
 \longrightarrow\frac1{1-M\omega^2/k}.
\]
这里分别使用 $\sin\beta/\beta\to1$、$\cos\beta\to1$，要求 $M\omega^2<k$。轻弹簧端球方程 $k(L-L_0)=M\omega^2L$也给 $(k-M\omega^2)L=kL_0$，恢复同一结果。

若 $\omega=0$，应直接解 $r''=0$。积分得 $r'(u)-r'(0)=\int_0^u0\dd s=0$，故 $r'=C$；再积分 $r(u)-r(0)=\int_0^u C\dd s=Cu$。固定端给 $r(0)=0$，末端零轴力给 $kL_0(C-1)=0$，所以 $C=1$、$L=L_0$，与零角速度的连续极限一致。
\end{drillsolution}
